\noindent\textbf{Scope and source control.}
This chapter covers every one of the 45 requested MCCME records: 229--271,
286 and 288.  The packet contains 99 scanned pages, including overlapping
talks and material outside the requested records.  Their actual handwritten
boundaries govern the entries.  In particular, P000358--P000359 continue
Perelomov's preceding-year talk, while P000450--P000451 and the neighbouring
Buchstaber pages in the 1973 packet are not assigned to year headings.
Record 288 is retained because it was explicitly requested, but its clear
handwritten date, 19 February 1973 on P000480, corrects the catalogue's 1972
date.  Its identifiers therefore begin \texttt{1972-288}; this inclusion does
not constitute a complete companion for 1972--73.  The catalogue is cited as
\citet{mccme2026gelfand}; all substantive status statements are dated November
2026.

\paragraph{The year's mathematical lines.}
The autumn sequence moves from Kirillov's prequantization and polarizations to
Springer fibres, finite incidence geometry and microlocal analysis.  Bernstein's
Godbillon--Vey exposition and Atiyah's spectral-flow discussion connect
characteristic classes, foliations and index theory.  The Vinogradov lectures
organize differential operators through modules, jets and their universal
properties; the formal-geometry discussion links infinite jets with
cohomology of Lie algebras of vector fields.

The spring has a sustained algebraic \(K\)-theory thread: quadratic and
hermitian forms, elementary matrices and Steinberg relations, polynomial and
Laurent constructions, Vaserstein's comparison programme, Segal's categories,
and Quillen's finite-field calculation.  Parallel sequences on polarized
Hodge structures and their variation, sites and derived categories supply the geometric
and categorical language.  Scattering theory returns through trace-class
perturbations and stationary constructions.  Morava's formal-group picture
anticipates the height-based organization of complex bordism.  Dyson's
identity and the continued formal-group discussion close the main packet;
the requested Dobrushin outlier is only a brief field-theory report.

\paragraph{What became of the material.}
The occurrence index has 262 items.  Counts are occurrences, not numbers
of independent discoveries; splitting definitions, methods and claims changes
their granularity.  The 4 catalogue headings and 1 reading
plan have no invented mathematical items.  The 257 core-year occurrences and 5
Dobrushin occurrences are distinguished in the report.

\begin{center}
\begin{tabular}{lrrrrrrr}
\toprule
Class & A & B & C & D & E & F & Total\\
\midrule
Occurrences & 34 & 18 & 3 & 0 & 182 & 25 & 262\\
\bottomrule
\end{tabular}
\end{center}

The predominance of class E reflects serious exposition rather than
unrecorded novelty.  Publication and first circulation must be separated:
Quillen's finite-field calculation, the Adams-operation fibre and the
odd-prime cohomology formula were already announced at the 1970 Congress
\citep[\S\S2--3]{quillen1971cohomology}, although the detailed Annals paper
appeared in 1972.  Conversely, the spectral-flow programme is followed into
\citet[\S7]{atiyahpatodisinger1976spectral3}.  Morava's early circulated notes
are identified as unpublished; later published constructions are not
backdated to them.

\paragraph{Questions and critical qualifications.}
The 9 problem-index entries preserve definite questions even when
they were answered quickly.  The elementary-relations question is linked to
Vaserstein's answer later in the same seminar sequence.  Dyson's exponent-26
question has later root-system explanations.  Pt\'ak's rational-\(p\)
uncertainty was already overtaken by a 1971 paper before his talk; the exact
general fate of the minimum-over-all-norms and extremal questions was not
established by the searches.  Such a verification gap is stated expressly,
and is not evidence that a problem remains open in 2026.  The trace-class
equivalence question has a dedicated later classification paper, but its
full criterion is not quoted without a verified theorem text.

No class C or D item supplies both a securely assessed present frontier and
a justified bounded research contribution in this run.  Accordingly, the
development section contains editorial verifications: the parity obstruction
to nonzero spectral flow for the full de Rham family, three exact algorithms
for Dyson's coefficients through 100, and an exact recovery of the classical
rank-32 lattice bound.  Independent re-derivations and checks found no error;
finite computation is never promoted to an all-degree theorem.

\paragraph{Continuity and limits.}
Earlier links include prequantization (\texttt{1969-178.3--4}),
the distribution-to-foliation problem (\texttt{1970-206.1--2},
\texttt{GP-1970-26}), and the
site and sheaf expositions (\texttt{1968-151}, \texttt{1970-184},
\texttt{1970-226}).  A thematic connection does not by itself justify merging
problem identifiers.  The 127 correction rows distinguish errors
of transcription or translation, uncertain readings, catalogue provenance,
and mathematical errors faithfully preserved in the notes.  Neither a
class label nor a verified bibliographic URL substitutes for the explicit
publication-match and verification fields attached to each occurrence.
\clearpage
